\section{EPS case studies}
\label{section:eps-case-study-results}

The EPS case studies verify the calculations required by FR-12 and FR-14 to FR-20 using BIRDS-4 and OreSat data.
Mission publications supply component ratings and selected calculation results \citep{birds2021batteryverification,shrestha2022birds4powerbudget,greenberg2021oresat,oresat2025solardesign}.
Explicit assumptions supply the remaining inputs.

\subsection{Verification setup}

The BIRDS electrical case uses a two-cell 3G30A solar section and a 3S2P NiMH battery.
The OreSat case uses a two-cell XTE-SF section and two independent 2S1P lithium-ion packs.
Both baseline sections face the Sun.
The comparisons also cover changes in illumination, cell technology, series and parallel arrangements, bus voltage and power-transfer method.
Appendix~\ref{appendix:eps-case-studies} groups the inputs and selected results by the requirement they verify.

Independent reference values were calculated directly from the input data, without calling the framework's numerical functions.
The reference summed component masses and volumes, calculated duty-weighted state powers, and obtained phase energies from mean power and duration.
Loss-adjusted phase energies determined source-power demand.
Battery sizing used eclipse demand, the assumed cycle-life/DoD relation and the energy margin.
Power and energy ratios were rounded upward to obtain integer cell counts.
These sums, products and ratios used decimal arithmetic, with 40 significant digits for products and ratios.
A separate orbital reference solved Kepler's equation at uniformly spaced times and tested Earth shadow in Cartesian coordinates.
This provided an independent check of the framework's shadow-boundary and panel-projection calculations.

Each nonzero numerical output had to agree with its reference within 1\%.
Zero values, integer counts and discrete outcomes required exact agreement.
The appendix reports maximum deviations and representative comparisons, with repeated outputs and deliberately failing cases omitted from the detailed tables.

\subsection{Results and discussion}

The numerical comparisons met the 1\% criterion for every assessed requirement in both mission cases.
Integer counts and discrete outcomes also matched their references.
No unexpected mismatch was found.
Appendix~\ref{appendix:eps-case-studies} separates the evidence for mass accounting, sizing, component choices, electrical configurations, conversion, photovoltaic characteristics, transfer paths and protection.

The largest deviation was 0.331\% in the BIRDS photovoltaic power comparison for FR-18.
The framework calculated $5.123\times0.050=0.25615$~W, compared with the published 0.257~W label.
The next largest deviation was 0.315\% for FR-14, where the OreSat cell ratings gave 2.34738~W against a published two-cell value of 2.34~W.
Both comparisons use printed values with limited precision.
They check consistency with the publications at the stated tolerance, while the separately calculated references provide a more precise check of the arithmetic.

The largest deviation from an independent calculation was 0.000864\% in the BIRDS battery-energy requirement.
The framework returned 4243.934~J against 4243.970~J from the reference.
This difference follows from the finite time grid used to estimate eclipse duration in the orbital reference.
Doubling that grid changed eclipse duration by at most 0.00205\%.
The minimum integer cell counts remained unchanged.
The small energy difference therefore had no effect on the sizing decisions checked here.

These results support using the framework for EPS accounting, preliminary sizing and comparisons of entered component options.
The evidence includes independent calculations of both the component quantities and the orbital geometry that links illumination to the power budget.
For FR-15 and FR-16, the technology and arrangement comparisons reproduced the expected changes in power, voltage, energy and capacity.
For FR-17, FR-19 and FR-20, terminal currents, transfer-path powers and protection comparisons agreed with the separately evaluated electrical relations.
This connects the numerical agreement to the decisions the framework is intended to support.

The remaining uncertainty concerns specific input models.
The complete verification curves use three photovoltaic points, and the battery life fits are assumed.
The published BIRDS 1.575~W measured peak remains unreconstructed.
OreSat's partial solar sweep supplies power-product checks but lacks the endpoints needed for a complete characteristic.
The case studies consequently provide stronger evidence for correct calculation and comparison than for reproducing either mission's measured energy performance.
